\section{Modos de fallo de METR y límites de su validez}

\subsection{Por qué fallan las trayectorias largas}

\videofigure{figures/metr-failure-modes.jpg}{Los fallos comunes incluyen una planeación o selección de herramientas inadecuada, razonamiento erróneo, abandono prematuro y ciclos de acciones repetidas.}{00:23:18--00:26:18}

Los errores por lo general no aparecen de golpe en el último paso, sino que se acumulan desde decisiones tempranas: se elige la herramienta equivocada, no se verifican los productos intermedios, se interpreta mal el estado del entorno, y a partir de ahí se sigue invirtiendo tokens y tiempo en la rama equivocada.

\begin{warningbox}{Más context no sustituye la gestión de estado}
Meter todo el historial completo de vuelta al modelo no equivale a que el agent recuerde el estado clave. Las trayectorias largas necesitan un grafo de tareas explícito, hechos ya verificados, problemas sin resolver, resúmenes de las salidas de las herramientas y puntos de retroceso.
\end{warningbox}

\begin{knowledgebox}{El comportamiento cíclico es un fallo de la política de detención}
Que un agent repita la misma acción una y otra vez suele indicar que no registró en su estado que "la acción ya se intentó y no funcionó", ni activó una escalación, un cambio de plan o una terminación. Detectar estados repetidos es más eficaz que simplemente aumentar el número máximo de pasos.
\end{knowledgebox}

\subsection{Si las tareas realmente representan el trabajo real}

\videofigure{figures/metr-validation.jpg}{METR valida sus conclusiones sobre el time-horizon desde perspectivas como el desorden del entorno, el sesgo en la dificultad de las tareas y las revisiones internas.}{00:26:18--00:28:06}

La validación necesita revisar: si la tarea es demasiado limpia, si la calificación automática pasa por alto problemas de calidad, si la estimación del tiempo humano es confiable, si el modelo ya había visto la pregunta, y si los errores del entorno se contabilizaron erróneamente como errores de capacidad. El curso enfatiza que la credibilidad de la evaluación proviene de estas auditorías, no de una bonita línea de tendencia.

\begin{importantbox}{El benchmark también necesita un verifier}
Así como el agent es evaluado por un verifier, el benchmark mismo también debe ser verificado: la representatividad de las tareas, las reglas de calificación, la varianza entre ejecuciones repetidas y las líneas base humanas pueden ser todas fuentes de error sistemático.
\end{importantbox}

\subsection{Resumen de la sección}

METR es adecuado para responder preguntas sobre la amplitud de las capacidades, pero no puede responder por sí solo preguntas sobre la calidad del producto de trabajo, el valor económico y el contexto real. Analizar las trayectorias fallidas y la validez del benchmark es una parte necesaria para interpretar correctamente el time horizon.
